\begin{proof}
    Докажем первую часть утверждения. Для доказательства того, что из дифференцируемости \(f\) в 
    точке \(z_0\) следует условие Коши-Римана, воспользуемся определением производной через 
    \(\hat{\alpha}\). Пусть \(A = B + iC\), где \(B, C \in \R\). Тогда определение перезаписывается
    в следующее выражение:
    \begin{align}
        [u(x + \Delta x, y + \Delta y) &+ iv(x + \Delta x, y + \Delta y)] - [u(x, y) + iv(x, y)] = 
        \nonumber \\ &= (B + iC)(\Delta x + i\Delta y) + (\beta (\Delta x, \Delta y) + i \gamma 
        (\Delta x, \Delta y)) \sqrt{(\Delta x)^2 + (\Delta y) ^ 2}. \tag{1}
    \end{align}

    Отдельно напишем вещественную и комплексные части

    \[u(x + \Delta x, y + \Delta y) - u(x, y) = B \Delta x - C \Delta y + \beta(\Delta x, 
    \Delta y)\sqrt{(\Delta x)^2 + (\Delta y) ^ 2};\]
    \[v(x + \Delta x, y + \Delta y) - v(x, y) = C \Delta x + B \Delta y + \gamma(\Delta x, 
    \Delta y)\sqrt{(\Delta x)^2 + (\Delta y) ^ 2}.\]

    Из многомерного анализа следует, что \(u_x = B, u_y = -C, v_x = C, v_y = B\), а также то что
    \(u, v\) дифференцируемы в точке \((x, y)\).

    Теперь докажем в обратную сторону. Из условий Коши-Римана следуют записи, сделанные при 
    доказательстве прошлого пункта при разделении на вещественную и комплексную части. Сложив их
    получим определение дифференцируемости функции \(f(z)\) в \(z\).

    Доказательство второго пункта утверждения следует из доказательства первого, ведь производная
    равна \(A = B + iC = u_x + i v_x\).
\end{proof}

\begin{note}
    Как и в вещественном анализе, верны теоремы об арифметических свойствах дифференцируемых 
    функций и о суперпозиции дифференцируемых функций.
\end{note}

\begin{example}
    \(e^z\) дифференцируема всюду в \(\Co\). \(|z|^2\) дифференцируема лишь в нуле. 
\end{example}

\begin{definition}
    \(f(z)\) регулярна на непустом открытом множестве \(A\), если \(f(z)\) дифференцируема в любой 
    точке \(A\). \(f(z)\) регулярна на непустом множестве \(E\), если \(f(z)\) регулярна на
    некотором открытом множестве \(A\), таком что \(E \subset A\).
\end{definition}

\begin{note}
    Сумма, произведение и частное регулярных функций~--- регулярны.
\end{note}

\begin{definition}
    Целыми функциями будем называть функции регулярные на \(\Co\).
\end{definition}

\Section{Функциональные последовательности и ряды}

Пусть есть множество \(E \subset \Co\) и на \(E\) заданы \(f_n(z)\). Если каждая \(f_n : E \mapsto 
\Co\), то последовательность \({f_n}\) будем называть функциональной.

\begin{definition}
    \({f_n}\) сходится равномерно к некоторой функции \(f\) на \(E\), если \(\forall \epsilon > 0
    \exists N \in \N \colon \forall n \ge N \ \& \ \forall z \in E \mapsto |f_n(z) - f(z)| < 
    \epsilon\). Обозначается как \(f_n \mathrel{\mathop{\rightrightarrows}\limits_{E}} f\).
\end{definition}

\begin{definition}
    Сумму вида \(S(z) = \sum\limits_{n = 0}^{\infty} f_n\) будем называть функциональным рядом. Ряд 
    сходится равномерно, если сходятся его частичные суммы.
\end{definition}

\begin{statement}
    Пусть \(\forall n \in \N \mapsto f_n(z) \in C(E)\) и \(f_n 
    \mathrel{\mathop{\rightrightarrows}\limits_{E}} f\). Тогда \(f(z) \in C(E)\).
\end{statement}

\begin{anote}
    Здесь и далее я буду использовать обозначения для сокращения записи дифференцируемости функции
    и её непрерывности. Если функция непрерывна на множестве \(E\), то это будет обозначаться 
    \(f(z) \in C(E)\). Аналогично, для дифференцируемости будем использовать \(f(z) \in D(E)\).
    Вероятно, вы уже встречались с подобными обозначениями на семинарах.
\end{anote}

\begin{statement}
    \textbf{(Признак сходимости равномерной сходимости функциональных рядов).} Пусть \(\forall n \ge
     N \ \& \ \forall z \in E \mapsto |f_n(z)| < a_n\) и ряд \(\sum\limits_{n = 0}^{\infty} a_n\)
     сходится. Тогда функциональный ряд \(f_n\) сходится равномерно на множестве \(E\).
\end{statement}


\Section{Интеграл вдоль кривой}

Пусть кривая \(\gamma\) КНД, а также задана функция \(f(z) = u(x, y) + iv(x, y) \in C(\gamma)\).
Интегралом первого рода будем называть \(\int\limits_{\gamma} f(z) |dz| = \int\limits_\alpha^\beta
f(z(t)) |z'(t)|dt\), где \(z(t)\)~--- параметризация кривой \(\gamma\). Название обусловлено подобию
данного интеграла криволинейному интегралу первого рода из курса кратных интегралов и теории поля.
Действительно \(\int\limits_\alpha^\beta f(z(t)) |z'(t)|dt = \int\limits_\alpha^\beta [u(x(t), y(t))
+iv(x(t), y(t))] \sqrt{(x'(t))^2 + (y'(t))^2} dt = \int\limits_\alpha^\beta uds + i
\int\limits_\alpha^\beta vds\). Введём свойства интеграла первого рода.

\begin{properties}
    \item \(\int\limits_{\gamma} |dz| = l(\gamma)\)~--- длина кривой.
    \item \(\int\limits_{\gamma} (af + bg) |dz| = a \int\limits_{\gamma} f |dz| + b 
    \int\limits_{\gamma} g |dz|\)~--- свойство линейности.
    \item Пусть \(\gamma = \gamma_1 \gamma_2\). Тогда \(\int\limits_{\gamma} f |dz| = 
    \int\limits_{\gamma_1} f |dz| + \int\limits_{\gamma_2} f |dz|\).
    \item \(\int\limits_{\gamma} f |dz| = \int\limits_{\gamma^{-1}} f |dz|\)
\end{properties}

Теперь введём интеграл второго рода \(\int\limits_{\gamma} f(z) dz = \int\limits_\alpha^\beta
f(z(t)) z'(t)dt\). Введём его свойства.

\begin{properties}
    \item \(\int\limits_{\gamma} (af + bg) dz = a \int\limits_{\gamma} f dz + b 
    \int\limits_{\gamma} g dz\)~--- свойство линейности.
    \item Пусть \(\gamma = \gamma_1 \gamma_2\). Тогда \(\int\limits_{\gamma} f dz = 
    \int\limits_{\gamma_1} f dz + \int\limits_{\gamma_2} f dz\).
    \item \(\int\limits_{\gamma} f |dz| = -\int\limits_{\gamma^{-1}} f |dz|\)
\end{properties}

Выведем оценку для интегралов первого и второго рода функции \(f \in C(\gamma)\) (также можно 
называть функцию \(f\) непрерывной на своём носителе):
\[|\int\limits_{\gamma} f(z) dz| \le \int\limits_{\gamma} |f(z)| |dz| \le l(\gamma) \cdot 
\max\limits_{\gamma}{|f(z)|}\]

\begin{proof}
    Докажем первое неравенство. Обозначим \(I = \int\limits_{\gamma} f(z) dz\). При \(I = 0\)
    неравенство верно. При \(I \ne 0\) можно использовать запись \(I = |I|e^{i\theta}\). Отсюда
    \(|I| = Ie^{-i\theta} = \int\limits_{\gamma} e^{-i\theta} f(z) dz = 
    \int\limits_{\alpha}^{\beta} e^{-i\theta} f(z(t)) z'(t) dt = \int\limits_{\alpha}^{\beta} 
    \re{[\dots]} dt + i \int\limits_{\alpha}^{\beta} \im{[\dots]} dt\). Заметим, что интеграл во
    втором слагаемом обращается в нуль. Так как для любого комплексного числа его действительная 
    часть меньше или равна его модулю, то продолжая оценку получим: \(\int\limits_{\alpha}^{\beta} 
    \re{[\dots]} dt + i \int\limits_{\alpha}^{\beta} \im{[\dots]} dt \le 
    \int\limits_{\alpha}^{\beta} |f(z(t))| \cdot |z'(t)| dt = \int\limits_{\gamma} |f(z)| |dz|\).

    Второе неравенство доказывается тривиально исходя из того, что функция в каждой точке своего
    носителя меньше или равна максимуму функции на носителе.
\end{proof}

\begin{statement}
    Если \(\gamma\)~--- КНД кривая, \(\forall n \ f_n(z) \in C(\gamma)\) и \(f_n \rightrightarrows
    f\), то \(\int\limits_{\gamma} f_n(z) dz \rightarrow \int\limits_{\gamma} f(z) dz\)
\end{statement}

\begin{proof}
    Зафиксируем \(\epsilon > 0\). Из условий теоремы, для него \(\exists N \colon \ \forall n \ge N
    \rightarrow \max\limits_{\gamma}{|f_n(z) - f(z)| \le \epsilon}\), а из оценки выше следует, что
    \(|\int\limits_{\gamma} (f_n(z) - f(z)) dz| \le l(\gamma) \cdot
    \max\limits_{\gamma}{|f_n(z) - f(z)|} \le l(\gamma)\epsilon\).
\end{proof}

Аналогичное утверждение можно также сформулировать для рядов:

\begin{statement}
    Если \(\gamma\)~--- КНД кривая, \(\forall n \ f_n(z) \in C(\gamma)\) и 
    \(\sum\limits_{n=0}^{\infty}f_n \rightrightarrows f\), то \(\int\limits_{\gamma} 
    \sum\limits_{n=0}^{\infty} f_n(z) dz \rightarrow \int\limits_{\gamma} f(z) dz\)
\end{statement}

\begin{definition}
    Пусть задана КНД кривая \(\gamma\) с параметризацией \(z = z(t)\) на \([a, b]\) и регулярное на
    \(\gamma\) отображение \(w = f(z)\). \textit{Образом} \(\gamma\) будем называть кривую с 
    параметризацией \(w = f(z(t))\).
\end{definition}

\begin{statement}
    Пусть задана КНД кривая \(\gamma\) с параметризацией \(z = z(t)\) на \([a, b]\), регулярное на
    \(\gamma\) отображение \(w = f(z)\), а также \(F(u) \in C(\Gamma)\), где \(\Gamma = f(\gamma)\).
    Тогда \(\int\limits_{\Gamma}F(w)dw = \int\limits_{\gamma} F(f(z))\cdot f'(z) dz\).
\end{statement}

\begin{proof}
    \(w(t) = f(z(t))\). Расписав правую и левую части в пределах параметризации и воспользовавшись 
    теоремой о дифференцируемости сложной функции получим тождественные равенства.
\end{proof}